\documentclass{article}
\usepackage{amsmath,amssymb}
\usepackage[margin=1in]{geometry}
\usepackage[UTF8]{ctex}
\begin{document}

\begin{align}
\text{(MP)}\qquad \text{minimize}\quad &\sum_{t\in T}c_t x_t \notag\\
\text{subject to}\quad &\sum_{t\in T}a_{it}x_t=1, &&i\in C, \tag{2}\\
&x_t\in\{0,1\}, &&t\in T. \notag
\end{align}
对最小数量的旅行树所施加的打印强化措施如下：
\begin{equation}
\sum_{t\in T}x_t\geq
\left\lceil\frac{\sum_{i\in C}\mu_c}{\theta\bar Q}\right\rceil.
\tag{3}
\end{equation}
\noindent\textit{原文的索引问题。}等式（3）在由$i\in C$索引的求和项下方打印出$\mu_c$。客户指定的含义被定义为$\mu_i$、$i\in C$；在（3）中没有绑定索引$c$。打印出的表达式保持不变，无需修复。

\section*{固定的数据和操作假设}
该问题定义在一个无向加权图上。
\[
G=(V,E),\qquad V=D\cup S\cup C,
\]
其中，$D$、$S$和$C$分别代表了仓库、卫星站和客户的集合。边集被划分为两个层次化的边集。
\begin{align*}
E^1&=\{(i,j)\mid i\in D\cup S,\ j\in S,\ i\neq j\},\\
E^2&=\{(i,j)\mid i\in S\cup C,\ j\in C,\ i\neq j\}.
\end{align*}
第一梯队（FE）的车辆会经过$E^1$，而第二梯队（SE）的车辆则会经过$E^2$。对于每一个$e\in E$，固定的重量$c_e$就是穿越边$e$的运输成本。

每个梯队都拥有无限数量的同质车辆。各车辆的载货能力满足…
\[
Q^{\mathrm{FE}}>Q^{\mathrm{SE}}.
\]
在卫星上，货物从FE车辆中卸下后，直接装载到SE车辆中，无需进行中间存储。不同FE车辆之间的货物运输是不允许的。多个FE车辆可以访问同一颗卫星，但前提是它们为不同的SE车辆提供货物。

对于每一个$i\in C$，要求$\xi_i$是一个非负整数随机变量，其均值为$\mu_i$，方差为$\sigma_i^2$：
\[
\xi_i\in\{0,1,2,\ldots\},\qquad
\mathbb E[\xi_i]=\mu_i,\qquad
\operatorname{Var}(\xi_i)=\sigma_i^2,
\qquad i\in C.
\]
联合需求分布是固定数据。不假设任何独立性：不同客户的订单可能存在统计上的依赖性。固定的概率约束水平$\eta$指定了每个SE路线上的需求必须在其容量范围内。一般问题定义中并未对$\eta$施加额外的数值限制。

\section*{路线与可行的游览路线集合}
\paragraph{第一梯队路线（定义1）。}
FE路线是一系列的步骤或过程。
\[
f=(n_0,n_1,\ldots,n_m,n_{m+1})
\]
与…一起
\begin{align*}
m&\geq1, & n_0&=n_{m+1}\in D,\\
n_i&\in S &&\text{for all }i\in\{1,\ldots,m\},\\
n_i&\neq n_j &&\text{for all }i,j\in\{1,\ldots,m\},\ i\neq j.
\end{align*}

\paragraph{第二梯队路线（定义2）。}
有效的 SE 路径是一个序列。
\[
r=(n_0,n_1,\ldots,n_m,n_{m+1})
\]
与…一起
\begin{align*}
m&\geq1, & n_0&=n_{m+1}\in S,\\
n_i&\in C &&\text{for }i\in\{1,\ldots,m\},\\
n_i&\neq n_j &&\text{for all }i,j\in\{1,\ldots,m\},\ i\neq j,
\end{align*}
并令 $C'=\{n_1,\ldots,n_m\}$，则有
\begin{equation}
\mathbb P\left(\sum_{i\in C'}\xi_i\leq Q^{\mathrm{SE}}\right)\geq\eta.
\tag{1}
\end{equation}
约束条件（1）适用于可行旅行树中每条SE路线，而不适用于整个旅行树的随机需求总和。

\paragraph{Tour-tree（定义3及继承的可行性）。}
一个“tour-tree”由一条单一的FE路线以及该FE路线所提供的所有SE路线组成。其序列表示方式如下：
\[
t=(n_0,(n_1,R_1),\ldots,(n_m,R_m),n_{m+1}),
\]
其中
\[
m\geq1,\qquad n_0=n_{m+1}\in D,\qquad
n_k\in S\quad\text{for }k\in\{1,\ldots,m\},
\]
$R_k$是从卫星$n_k$出发的SE路线集合。 underlying FE路线遵循定义1，而每个$R_k$中的路线则遵循定义2。在基本游览树中，客户最多只能被访问一次；本文将这些基本游览树与后续的算法松弛方法区分开来，后者允许重复访问。

每条 SE 路线都会让其卫星处于满负荷状态。因此，每条 FE 路线所供应的资源最多为…
\[
\theta=\left\lfloor\frac{Q^{\mathrm{FE}}}{Q^{\mathrm{SE}}}\right\rfloor
\]
SE路线。在定义3的表示法中，这个条件就是
\[
\sum_{k=1}^{m}|R_k|\leq\theta
\qquad\text{for every feasible tour-tree }t.
\]
上述关于直接转让和单一FE供应商的条件同样适用。这些条件是用来定义各列的，并非MP中的额外决策。

\section*{列系数与决策变量}
$T$是指根据前述定义和运营假设所得到的所有可行旅行树集合。对于$t\in T$和$E(t)$，分别表示在旅行树$t$中经过的边集，而其固定的运输成本则保持不变。
\[
c_t=\sum_{e\in E(t)}c_e,\qquad t\in T.
\]
$E(t)$和$c_t$的集合式表达保持原样；没有添加边遍历重数的约定。

固定的二进制关联参数为
\[
a_{it}=\begin{cases}
1,&\text{if customer }i\text{ is visited in tour-tree }t,\\
0,&\text{otherwise},
\end{cases}
\qquad i\in C,\ t\in T.
\]
唯一的MP决策变量是$x_t$和$t\in T$：当其中一个的值为1时，就会选择tour-tree$t$。方程(2)确保每个客户都只被一个选定的tour-tree所覆盖。

\section*{(3)中定义的“bound”的定义}
第3节通过以下性质来定义$\bar Q$：
\[
\sum_{i\in C'}\mu_i>\bar Q
\quad\Longrightarrow\quad
\mathbb P\left(\sum_{i\in C'}\xi_i\leq Q^{\mathrm{SE}}\right)<\eta,
\qquad\text{for every }C'\subset C.
\]
所引用的定义可以在第5.3节中找到（物理PDF页面15-17）。对于客户组合，该定义中有所说明。
\[
\mu(C')=\sum_{i\in C'}\mu_i,\qquad C'\subset C.
\]
其上限可行性界限$\bar\mu$必须满足以下条件：
\[
\mu(C')>\bar\mu
\quad\Longrightarrow\quad
\mathbb P\left(\sum_{i\in C'}\xi_i\leq Q^{\mathrm{SE}}\right)<\eta,
\qquad\text{for every }C'\subset C.
\]
本文将式(3)中使用的固定参数设定为这个界限：
\[
\boxed{\bar Q=\bar\mu.}
\]
要获取$\bar\mu$，需要按照打印出的步骤进行初始化操作。
\[
\mu^*=\sum_{i\in C}\mu_i
\]
并且将$\mu^*$减少一，直到找到一种客户组合$C'\subset C$，使得…
\[
\mu(C')=\mu^*,\qquad
\mathbb P\left(\sum_{i\in C'}\xi_i\leq Q^{\mathrm{SE}}\right)\geq\eta;
\]
然后，它设置$\bar\mu=\mu^*$的值。这里的$\mu^*$是一个预处理搜索值，而不是一个MP决策变量。在这个公式中，没有使用任何情景选择优化问题、采样近似方法或基于方差的可行性测试来替代(1)。

\noindent\textit{印刷装订结构的适用范围。}单位减少量在算法2中有所说明，不过总体问题陈述并不要求使用整数表示法。这里没有添加任何四舍五入的规则或特殊情况下的终止规则。第5.3节中还包含了一句话，即除了$\bar\mu$之外，超过容量的概率超过$\eta$。不过，这句话并不是(1)式所定义的违规条件；上述条件仍然保留了(1)式被违反的明确要求，即需求不超过容量的概率严格小于$\eta$。

\paragraph{来源范围。}
该模型为 MP，其打印后的强化项为（3），其中包含了第1至第3条定义，以及第3节中的操作假设（物理PDF页面5-6）。此外，第5.3节中的定义$\bar Q=\bar\mu$也被包含在内（物理PDF页面15-17）。方差限制附录以及定价和可行性搜索子问题并未对该MP提出额外的约束条件。也没有对（1）进行确定性替换，或者放宽二进制变量的限制。

\end{document}
